\clearpage
\VerdierPhysicalPageAnchor{0057}%
\VerdierSemanticAnchor{ch1-1-5-7-continuation}
显现出五个竖直图，依次为：

\VerdierSemanticAnchor{ch1-1-5-7-vertical}
\[
  T'(\sigma')\Delta(\sigma,\tau),\quad
  T'(\sigma)\Delta(\sigma',\tau),\quad
  \Delta(\sigma,\tau)T(\sigma'),\quad
  \Delta(\sigma',\tau)T(\sigma),\quad
  \Delta(\sigma+\sigma',\tau).
\]

同样，它还显现出四个水平图，依次为：

\VerdierSemanticAnchor{ch1-1-5-7-horizontal}
\[
  D(\sigma,\sigma')T(\tau),\quad
  D(\sigma',\sigma)T(\tau),\quad
  T'(\tau)D(\sigma,\sigma'),\quad
  T'(\tau)D(\sigma',\sigma).
\]

\VerdierSemanticAnchor{ch1-1-5-7-conclusion}
由归纳假设，图 $D(\sigma,\sigma')$ 与 $D(\sigma',\sigma)$ 分别为
$\alpha(\sigma,\sigma')$-交换和 $\alpha(\sigma',\sigma)$-交换的。
考察三种可能情形，便得到断言 (1.5.4.3) 的第一部分。交换 $\sigma$ 与
$\tau$ 的角色即可得到其第二部分。在一般情形中，需要引入的图更复杂；
但利用恒等式 (1.2.1.3)，论证可归结为上面给出的论证。

\VerdierSemanticAnchor{ch1-1-5-8}
\noindent\textbf{注 1.5.8.}\quad
命题 (1.5.2) 特别适用于余循环 $\alpha$ 平凡的情形。这时可以验证，
即使范畴 $\mathcal C'$ 不再必为加法范畴，该命题及其证明仍可逐字照搬。

\VerdierSemanticAnchor{ch1-1-5-9}
\noindent\textbf{1.5.9.}\quad
设

\VerdierSemanticAnchor{ch1-1-5-9-functors}
\[
  (F_1,p_1),(F_2,p_2)\colon
  (\mathcal C,\mathbb Z^n,T,c)
  \longrightarrow
  (\mathcal C',\mathbb Z^n,T',c')
\]
为分别由两个正规化余循环
$\alpha_1,\alpha_2\in Z^2(\mathbb Z^n,\mathbb Z/2\mathbb Z)$
扭曲的两个 $\mathbb Z^n$-函子，并设
$\beta\colon\mathbb Z^n\to\mathbb Z/2\mathbb Z$ 为满足

\VerdierSemanticAnchor{ch1-1-5-9-eq1}
\begin{equation}
  \alpha_2(\sigma_1,\sigma_2)\alpha_1^{-1}(\sigma_1,\sigma_2)
  =\beta(\sigma_2)\beta(\sigma_1\sigma_2)^{-1}\beta(\sigma_1),
  \qquad \sigma_1,\sigma_2\in\mathbb Z^n,
  \tag{1.5.9.1}
\end{equation}
的映射，又设 $u\colon(F_1,p_1)\to(F_2,p_2)$ 为函子的一个
$\beta$-态射 (1.4.6)。令 $e_1,\ldots,e_n$ 为 $\mathbb Z^n$ 的典范基。
对每个 $i\in[1,n]$，下图为 $\beta(e_i)$-交换的：

\VerdierSemanticAnchor{ch1-1-5-9-eq2}
\begin{equation}
  \resizebox{0.78\textwidth}{!}{%
    $\begin{tikzcd}[ampersand replacement=\&,column sep=6em,row sep=4em]
      F_1T(e_i)
        \arrow[r,"{p_1(e_i)}"]
        \arrow[d,"{u*T(e_i)}"']
        \& T'(e_i)F_1
          \arrow[d,"{T'(e_i)*u}"] \\
      F_2T(e_i)
        \arrow[r,"{p_2(e_i)}"]
        \& T'(e_i)F_2
    \end{tikzcd}$%
  }\;.
  \tag{1.5.9.2}
\end{equation}
